\section*{Capítulo \ref{ch:g11:fundamental-interactions} --- Las interacciones fundamentales}

\begin{solution}{exo:g11:fundamental-interactions:1}
Protón: $10^{-15}$ (planta del \omterm{def:g11:fundamental-interactions:ladder}{nucleón}); \omterm{def:g11:fundamental-interactions:ladder}{molécula} de agua: $10^{-10}$
(planta de la \omterm{def:g11:fundamental-interactions:ladder}{molécula} --- una \omterm{def:g11:fundamental-interactions:ladder}{molécula} suelta tiene tamaño atómico);
arena: $10^{-4}$ (materia cotidiana); Tierra: $10^{7}$ (planeta); Vía
Láctea: $10^{21}$ (galaxia). Todo en \omterm{def:g10:orders-of-magnitude:unit}{metros}.
\end{solution}

\begin{solution}{exo:g11:fundamental-interactions:2}
Carga negativa: ha \emph{ganado}
$N = \num{4.8e-9}/\num{1.6e-19} = \num{3.0e10}$ electrones. El pelo
queda con $+\qty{4.8e-9}{C}$, por la conservación de la carga
(\cref{prop:g11:fundamental-interactions:conservation}).
\end{solution}

\begin{solution}{exo:g11:fundamental-interactions:3}
$F = \num{8.99e9} \times \dfrac{\num{2.0e-6} \times \num{3.0e-6}}{0.30^2}
= \dfrac{\num{5.39e-2}}{0.090} \approx \qty{0.60}{N}$, atractiva (signos
distintos). Las dos \omterm{def:g10:inertia:force}{fuerzas} tienen el mismo módulo y sentidos opuestos.
\end{solution}

\begin{solution}{exo:g11:fundamental-interactions:4}
$F_0/4$; \ $9F_0$; \ $2F_0$; \ $F_0$ (numerador $\times 4$, denominador
$\times 4$).
\end{solution}

\begin{solution}{exo:g11:fundamental-interactions:5}
(a) la \omterm{def:g11:fundamental-interactions:four}{gravitación}; (b) la electromagnética; (c) la fuerte; (d) la
débil; (e) la electromagnética.
\end{solution}

\begin{solution}{exo:g11:fundamental-interactions:6}
$F_E = \num{8.99e9} \times (\num{1.60e-19})^2/(\num{2.0e-15})^2
= \num{2.30e-28}/\num{4.0e-30} \approx \qty{58}{N}$;
$F_G = \num{6.67e-11} \times (\num{1.67e-27})^2/\num{4.0e-30}
\approx \qty{4.7e-35}{N}$. Cociente $F_E/F_G \approx \num{1.2e36}$. Lo
que sostiene el \omterm{def:g11:fundamental-interactions:ladder}{núcleo} es la \omterm{def:g11:fundamental-interactions:strong}{interacción fuerte}; la gravedad es
irrelevante.
\end{solution}

\begin{solution}{exo:g11:fundamental-interactions:7}
$F_E = \num{2.30e-28}/(\num{5.3e-11})^2 = \num{2.30e-28}/\num{2.81e-21}
\approx \qty{8.2e-8}{N}$;
$F_G = \num{6.67e-11} \times \num{1.67e-27} \times \num{9.11e-31} /
\num{2.81e-21} \approx \qty{3.6e-47}{N}$. Cociente
$\approx \num{2.3e39}$: lo que mantiene unidos a los \omterm{def:g11:fundamental-interactions:ladder}{átomos} es la
\omterm{def:g11:fundamental-interactions:four}{interacción electromagnética}.
\end{solution}

\begin{solution}{exo:g11:fundamental-interactions:8}
Por contacto la carga se reparte: $q_A = q_B = \qty{+4.0e-9}{C}$ (el
total sigue siendo $\qty{8.0e-9}{C}$). Entonces
$F = \num{8.99e9} \times (\num{4.0e-9})^2/0.10^2 = \qty{1.4e-5}{N}$,
repulsiva (signos iguales).
\end{solution}

\begin{solution}{exo:g11:fundamental-interactions:9}
$d = \sqrt{k q^2/F} = \sqrt{\num{8.99e9} \times \num{1.0e-12}/0.90}
\approx \qty{0.10}{m}$. Una \omterm{def:g10:inertia:force}{fuerza} $9$ veces menor exige una distancia
$3$ veces mayor: $d = \qty{0.30}{m}$.
\end{solution}

\begin{solution}{exo:g11:fundamental-interactions:10}
$F = \num{2.30e-28}/(\num{1.0e-10})^2 = \qty{2.3e-8}{N}$. La \omterm{def:g11:fundamental-interactions:strong}{interacción
fuerte} no aporta nada: $10^{-10}$~m es $10^{5}$ veces su \omterm{def:g11:fundamental-interactions:strong}{alcance}. Las
\omterm{def:g11:fundamental-interactions:ladder}{moléculas} están ligadas por la \omterm{def:g11:fundamental-interactions:four}{interacción electromagnética}.
\end{solution}

\begin{solution}{exo:g11:fundamental-interactions:11}
$q = \num{e23}/\num{e6} \times \num{1.6e-19} = \qty{1.6e-2}{C}$ en cada
moneda. $F = \num{8.99e9} \times (\num{1.6e-2})^2/1.0^2 \approx
\qty{2.3e6}{N}$ --- el \omterm{def:g10:universal-gravitation:weight}{peso} de
$\num{2.3e6}/9.81 \approx \qty{2.3e5}{kg}$, un tren de mercancías
cargado, a partir de un millonésimo de los electrones. La materia
cotidiana tiene que ser (y es) neutra con mucha más precisión que una
parte en un millón.
\end{solution}

\begin{solution}{exo:g11:fundamental-interactions:12}
$F = \num{2.30e-28}/(\num{1.4e-14})^2 = \num{2.30e-28}/\num{1.96e-28}
\approx \qty{1.2}{N}$ --- enorme todavía a esta escala. No: están a $14$
veces el \omterm{def:g11:fundamental-interactions:strong}{alcance} fuerte; cada protón solo está pegado a sus vecinos más
próximos. Parejas: $\frac{92 \times 91}{2} = 4186$, y \emph{todas} se
repelen a cualquier distancia, mientras que el pegamento no crece con el
tamaño. Los neutrones añaden atracción sin repulsión y separan a los
protones; pero la repulsión crece con el cuadrado del número de protones
y el pegamento solo con el número de vecinos, así que más allá de unos
$80$ protones ningún presupuesto de neutrones cuadra las cuentas ---
esos \omterm{def:g11:fundamental-interactions:ladder}{núcleos} se desintegran
(\cref{ch:g11:nucleus-radioactivity}).
\end{solution}

\begin{solution}{exo:g11:fundamental-interactions:13}
Equilibrio: en vertical $T\cos\theta = mg$ y en horizontal
$T\sin\theta = F$, luego $F = mg\tan\theta$. Aquí
$\sin\theta = 3.0/50 = 0.060$ y, para ángulos pequeños,
$\tan\theta \approx \sin\theta = 0.060$. Entonces
$F = \num{1.0e-3} \times 9.81 \times 0.060 \approx \qty{5.9e-4}{N}$;
$q = d\sqrt{F/k} = 0.060 \times \sqrt{\num{5.9e-4}/\num{8.99e9}}
\approx \qty{1.54e-8}{C}$;
$N = \num{1.54e-8}/\num{1.6e-19} \approx \num{9.6e10}$ \omterm{def:g11:fundamental-interactions:elementary}{cargas
elementales}.
\end{solution}

\begin{solution}{exo:g11:fundamental-interactions:14}
$F = \num{6.67e-11} \times \num{5.97e24} \times \num{7.35e22} /
(\num{3.84e8})^2 = \num{2.93e37}/\num{1.47e17} \approx \qty{2.0e20}{N}$.
Imponiendo $kQ^2/d^2 = F$:
$Q = d\sqrt{F/k} = \num{3.84e8} \times \sqrt{\num{2.0e20}/\num{8.99e9}}
\approx \qty{5.7e13}{C}$. Eso son
$N = \num{5.7e13}/\num{1.6e-19} \approx \num{3.6e32}$ electrones, de
masa $\num{3.6e32} \times \num{9.11e-31} \approx \qty{3.2e2}{kg}$: unos
pocos cientos de \omterm{def:g10:orders-of-magnitude:unit}{kilogramos} de electrones podrían hacer el trabajo de
los $\qty{7.35e22}{kg}$ de la Luna --- la \omterm{def:g10:inertia:force}{fuerza} de la electricidad y la
debilidad de la gravedad en un solo número.
\end{solution}

\begin{solution}{exo:g11:fundamental-interactions:15}
$F = \num{2.30e-28}/(\num{2.8e-10})^2 = \num{2.30e-28}/\num{7.84e-20}
\approx \qty{2.9e-9}{N}$; \omterm{def:g10:universal-gravitation:weight}{peso}
$P = \num{5.9e-26} \times 9.81 \approx \qty{5.8e-25}{N}$; cociente
$F/P \approx \num{5e15}$. Cada enlace tiene una intensidad
\emph{fija}, pero el \omterm{def:g10:universal-gravitation:weight}{peso} crece con toda la masa que hay encima: si se
apila roca, la carga sobre la capa inferior crece sin límite y sus
enlaces no. Hacia los \qty{e4}{m} de roca, la carga aplasta los enlaces
eléctricos --- la gravedad gana no por su intensidad por partícula, sino
porque \emph{se suma} cuando la carga se cancela.
\end{solution}

\begin{solution}{pb:g11:fundamental-interactions:1}
\textbf{1.} \omterm{def:g11:fundamental-interactions:ladder}{Quark} $< \qty{e-18}{m}$; \omterm{def:g11:fundamental-interactions:ladder}{nucleón} $\qty{e-15}{m}$; \omterm{def:g11:fundamental-interactions:ladder}{núcleo} de
$10^{-15}$ a $\qty{e-14}{m}$; \omterm{def:g11:fundamental-interactions:ladder}{átomo} $\qty{e-10}{m}$; \omterm{def:g11:fundamental-interactions:ladder}{molécula}
$\qty{e-9}{m}$; objeto cotidiano $\qty{1}{m}$; planeta $\qty{e7}{m}$;
sistema planetario $\qty{e11}{m}$; galaxia $\qty{e21}{m}$.

\textbf{2.} Factor de escala $10^{5}$: el \omterm{def:g11:fundamental-interactions:ladder}{átomo} mide
$\qty{1}{cm} \times 10^{5} = \qty{1}{km}$.

\textbf{3.} $(10^{-15}/10^{-10})^3 = 10^{-15}$: la materia está vacía en
un $\num{99.9999999999999}\,\%$.

\textbf{4.} $\num{7e-5}/\num{e-10} = \num{7e5}$ \omterm{def:g11:fundamental-interactions:ladder}{átomos} --- alrededor de
un millón a lo ancho de un cabello.

\textbf{5.} De $10^{-18}$ a $10^{21}$: $39$ potencias de diez.

\textbf{6.} $F_E = \num{8.99e9} \times (\num{1.60e-19})^2 /
(\num{5.3e-11})^2 \approx \qty{8.2e-8}{N}$.

\textbf{7.} $F_G = \num{6.67e-11} \times \num{1.67e-27} \times
\num{9.11e-31}/(\num{5.3e-11})^2 \approx \qty{3.6e-47}{N}$; cociente
$F_E/F_G \approx \num{2.3e39}$. La electricidad sostiene el \omterm{def:g11:fundamental-interactions:ladder}{átomo}; la
gravedad es una espectadora.

\textbf{8.} Sin gravedad: no cambia nada medible (está $10^{39}$ por
debajo). Sin electricidad: ni \omterm{def:g11:fundamental-interactions:ladder}{átomos}, ni \omterm{def:g11:fundamental-interactions:ladder}{moléculas}, ni materia alguna.

\textbf{9.} El ``contacto'' es la repulsión eléctrica entre las nubes de
electrones de los \omterm{def:g11:fundamental-interactions:ladder}{átomos} cuando se acercan a unos $\qty{e-10}{m}$: nada
se toca nunca. La reacción del suelo y el \omterm{def:g10:inertia:inventory}{rozamiento} son la ley de
Coulomb al por mayor.

\textbf{10.} $q = \num{e23} \times \num{e-9} \times \num{1.6e-19} =
\qty{1.6e-5}{C}$ en cada una;
$F = \num{8.99e9} \times (\num{1.6e-5})^2 \approx \qty{2.3}{N}$ --- una
\omterm{def:g10:inertia:force}{fuerza} perfectamente visible a partir de un electrón de cada mil
millones. La materia cotidiana es neutra con mucha más precisión que
$10^{-9}$.

\textbf{11.} La \omterm{def:g11:fundamental-interactions:four}{interacción electromagnética}. La \omterm{def:g10:inertia:force}{fuerza} fuerte queda
fuera de \omterm{def:g11:fundamental-interactions:strong}{alcance} más allá de unos $\qty{e-14}{m}$ --- $10^{4}$ veces
menos que un solo \omterm{def:g11:fundamental-interactions:ladder}{átomo}.

\textbf{12.} $F = \num{8.99e9} \times (\num{1.60e-19})^2 /
(\num{1.0e-15})^2 \approx \qty{2.3e2}{N}$.

\textbf{13.} $a = F/m_p = 230/\num{1.67e-27} \approx
\qty{1.4e29}{m/s^2}$, unas $\num{1.4e28}$ veces $g$: ninguna \omterm{def:g10:inertia:force}{fuerza}
cotidiana se le acerca.

\textbf{14.} Más intensa que \qty{230}{N} a $\qty{e-15}{m}$ (los \omterm{def:g11:fundamental-interactions:ladder}{núcleos}
existen); actuando por igual sobre protones y neutrones, ciega a la
carga (los neutrones también quedan agarrados); y de \omterm{def:g11:fundamental-interactions:strong}{alcance} en torno a
$\qty{e-15}{m}$ (los \omterm{def:g11:fundamental-interactions:ladder}{núcleos} siguen siendo diminutos --- el pegamento no
llega más lejos).

\textbf{15.} $F = \num{2.30e-28}/(\num{1.4e-14})^2 \approx
\qty{1.2}{N}$. No: están a $14$ veces el \omterm{def:g11:fundamental-interactions:strong}{alcance}. La repulsión actúa
sobre \emph{todas} las $\frac{92\times91}{2} = 4186$ parejas de
protones; el pegamento solo liga a los vecinos más próximos, así que
crecer favorece a la repulsión. Los neutrones añaden pegamento y
separación sin añadir repulsión --- hasta que, más allá de unos $80$
protones, ninguna mezcla cuadra y el \omterm{def:g11:fundamental-interactions:ladder}{núcleo} se desintegra.

\textbf{16.} La \omterm{def:g11:fundamental-interactions:weak}{interacción débil} convierte un neutrón en protón (y al
revés), lo que produce la desintegración $\beta$ --- transforma; no liga
nunca.

\textbf{17.} $F = \num{6.67e-11} \times \num{5.97e24} \times
\num{7.35e22}/(\num{3.84e8})^2 \approx \qty{2.0e20}{N}$.

\textbf{18.} $Q_E = \num{1.8e51} \times \num{1.6e-19} =
\qty{2.9e32}{C}$ y $Q_M = \qty{3.5e30}{C}$. Igualar las \omterm{def:g10:inertia:force}{fuerzas} exige
$f^2\,k\,Q_E Q_M = G M_E M_M$:
\[
f = \sqrt{\frac{\num{2.93e37}}{\num{8.99e9} \times \num{2.9e32}
\times \num{3.5e30}}} = \sqrt{\num{3.2e-36}} \approx \num{1.8e-18} :
\]
un desequilibrio de \emph{dos partes en $10^{18}$} bastaría para
cancelar la gravedad.

\textbf{19.} La carga viene en dos signos y se cancela (la pregunta 18
muestra con qué perfección), de modo que la \omterm{def:g10:inertia:force}{fuerza} más intensa se apaga
a sí misma a gran escala; la masa tiene un solo signo, se suma para
siempre y nada la apantalla.

\textbf{20.} \omterm{def:g11:fundamental-interactions:ladder}{Núcleo}: la \omterm{def:g11:fundamental-interactions:strong}{interacción fuerte} puede más que la repulsión de
los protones a $\qty{e-15}{m}$. Del \omterm{def:g11:fundamental-interactions:ladder}{átomo} a la montaña: la \omterm{def:g11:fundamental-interactions:four}{interacción
electromagnética} liga los electrones a los \omterm{def:g11:fundamental-interactions:ladder}{núcleos} y los \omterm{def:g11:fundamental-interactions:ladder}{átomos} entre
sí. Planeta y más allá: la \omterm{def:g11:fundamental-interactions:four}{gravitación}, sin oposición una vez que la
carga se cancela, sostiene planetas, \omterm{def:g10:universal-gravitation:satellite}{órbitas} y galaxias. El mundo no se
derrumba porque cada planta está apuntalada por un ligante más rígido
que la carga, y no se desintegra porque la materia es neutra hasta unas
$10^{-18}$ partes --- así que cada escala tiene exactamente una \omterm{def:g10:inertia:force}{fuerza}
al mando, y es atractiva donde tiene que serlo.
\end{solution}
